% =====================================================================
%  Appendices — referenced from the main text. Place these AFTER
%  \appendix in your paper, or in a separate appendix file.
% =====================================================================

\appendix

% ---------------------------------------------------------------------
\section{Pseudocode: sample-once PATE feature aggregation}
\label{app:pseudocode}

Algorithm~\ref{alg:precompute} gives the precompute routine that
implements the sample-once-per-image release with the PATE
multi-teacher aggregation mechanism (\S\ref{ssec:threat},
\S\ref{ssec:pate}, \S\ref{ssec:impl}). It is invoked once before the
student training loop. Throughout, $\phi_T^{(k)}$ denotes teacher $k$'s
(frozen) encoder trained on cohort $\mathcal{D}_k$,
$\mathrm{clip}(\cdot;\widehat{c}^{(k)})$ per-channel $L_2$ clipping to
caps $\widehat{c}^{(k)}$, and $\eta\sim\mathcal{N}(0,\sigma^2)$ draws
independent Gaussian noise per channel and spatial location.

\begin{algorithm}[t]
\caption{Sample-once precompute of PATE-aggregated noisy bottlenecks.}
\label{alg:precompute}
\begin{algorithmic}[1]
\Require $K$ teachers $\{\phi_T^{(k)}\}_{k=1}^K$ trained on disjoint
         cohorts; private training images $\{x_i\}_{i=1}^{N}$;
         per-teacher per-channel caps $\widehat{c}^{(k)} \in \mathbb{R}^{C}$
         (computed once on public proxy); budget $\rho_{\mathrm{rel}}$.
\Ensure  noisy bottleneck cache
         $\mathcal{C}:\mathrm{img\_id}\!\to\!\mathbb{R}^{C\times H\times W}$
\State $\Delta \gets 2/K$                                            \Comment{Theorem~\ref{thm:pate-sens}}
\State $\sigma \gets \sqrt{C\,\Delta^{2} / (2\rho_{\mathrm{rel}})}$    \Comment{uniform allocation on aggregate}
\State $\mathcal{C} \gets \emptyset$
\For{$i = 1,\dots,N$}
  \For{$k = 1,\dots,K$}
    \State $z_i^{(k)} \gets \phi_T^{(k)}(x_i)$                         \Comment{teacher $k$ forward}
    \State $z_i^{\prime (k)} \gets \mathrm{clip}(z_i^{(k)}; \widehat{c}^{(k)}) / \widehat{c}^{(k)}$
                                                                       \Comment{per-channel norm $\le 1$}
  \EndFor
  \State $\bar z_i \gets \tfrac{1}{K}\sum_{k=1}^{K} z_i^{\prime (k)}$  \Comment{PATE aggregate}
  \State sample $\eta_i \in \mathbb{R}^{C\times H\times W}$ with
         $\eta_i^{(c)} \!\sim\! \mathcal{N}(0,\sigma^{2})$
  \State $\widetilde{z}_i \gets (\bar z_i + \eta_i)\odot \widehat{c}^{\mathrm{ref}}$
                                                                       \Comment{denormalise w/ ref cap}
  \State $\mathcal{C}[\mathrm{img\_id}(x_i)] \gets \widetilde{z}_i$
\EndFor
\State \Return $\mathcal{C}$
\end{algorithmic}
\end{algorithm}

\noindent During student training, the inner loop's teacher path is
replaced by a single look-up $\widetilde{z}_i \leftarrow
\mathcal{C}[\mathrm{img\_id}(x_i)]$, eliminating both per-iteration
teacher queries and per-iteration noise samples. The reference cap
$\widehat{c}^{\mathrm{ref}}$ in line~10 of
Algorithm~\ref{alg:precompute} can be any per-channel scaling
consistent across the dataset (we use the mean of $\widehat{c}^{(k)}$
across $k$); the choice is post-processing and does not affect the
DP guarantee.

% ---------------------------------------------------------------------
\section{Full proof of Theorem~\ref{thm:wf} (CANAL closed form)}
\label{app:proof}

We restate the optimisation problem from \eqref{eq:opt} for
convenience:
\begin{equation*}
\min_{\sigma_1,\ldots,\sigma_C>0}\; \sum_{c=1}^{C} s_c\,\sigma_c^{2}
\quad\text{s.t.}\quad
\sum_{c=1}^{C} \frac{\Delta_c^{2}}{2\,\sigma_c^{2}} \le \rho_{\mathrm{rel}}.
\end{equation*}

\paragraph{Change of variables.}
Let $u_c = 1/\sigma_c^{2}$, so $u_c > 0$ and $\sigma_c^{2} = 1/u_c$.
The problem becomes
\begin{equation}
\min_{u_1,\ldots,u_C>0}\; \sum_{c=1}^{C} \frac{s_c}{u_c}
\quad\text{s.t.}\quad
\sum_{c=1}^{C} \frac{\Delta_c^{2}}{2}\,u_c \le \rho_{\mathrm{rel}}.
\label{eq:dual}
\end{equation}
Both the objective $\sum_c s_c/u_c$ (convex in $u_c$ on the positive
orthant) and the linear constraint give a convex program with a
unique minimiser. The constraint is active at the optimum because the
objective is strictly decreasing in each $u_c$.

\paragraph{Lagrangian.}
Form the Lagrangian with multiplier $\lambda > 0$:
\begin{equation*}
\mathcal{L}(u,\lambda)
 \;=\; \sum_{c=1}^{C} \frac{s_c}{u_c}
       + \lambda\left(\sum_{c=1}^{C}\frac{\Delta_c^{2}}{2}u_c - \rho_{\mathrm{rel}}\right).
\end{equation*}
Stationarity $\partial\mathcal{L}/\partial u_c = 0$ gives
\begin{equation*}
-\frac{s_c}{u_c^{2}} + \frac{\lambda\,\Delta_c^{2}}{2} = 0
\;\;\Longrightarrow\;\;
u_c^{\star} \;=\; \sqrt{\frac{2\,s_c}{\lambda\,\Delta_c^{2}}}
\;=\; \sqrt{\frac{2}{\lambda}}\;\frac{\sqrt{s_c}}{\Delta_c}.
\end{equation*}

\paragraph{Solving for $\lambda$ via the active constraint.}
Substitute back into $\sum_c \tfrac{\Delta_c^{2}}{2}u_c^{\star} =
\rho_{\mathrm{rel}}$ and solve:
\begin{equation*}
\lambda \;=\; \frac{1}{2\,\rho_{\mathrm{rel}}^{2}}
              \Bigl(\sum_{c=1}^{C}\Delta_c\sqrt{s_c}\Bigr)^{2}.
\end{equation*}

\paragraph{Optimal $\sigma_c^{\star}$.}
Substituting $\lambda$ back into $u_c^{\star}$ and using
$\sigma_c^{\star} = 1/\sqrt{u_c^{\star}}$:
\begin{equation*}
\sigma_c^{\star}
 \;=\;
 \underbrace{\sqrt{\frac{1}{2\,\rho_{\mathrm{rel}}}\sum_{c=1}^{C}\Delta_c\sqrt{s_c}}}_{\displaystyle\kappa}
 \;\frac{\sqrt{\Delta_c}}{s_c^{1/4}},
\end{equation*}
which is exactly \eqref{eq:wf}.

\paragraph{Optimal objective and $R^{1/4}$ bound.}
The optimal value of the original objective is
\begin{equation}
J_{\mathrm{WF}}^{\star}
 \;=\; \frac{1}{2\,\rho_{\mathrm{rel}}}
       \Bigl(\sum_{c=1}^{C}\Delta_c\sqrt{s_c}\Bigr)^{2}.
\label{eq:wf-obj}
\end{equation}
For uniform $\sigma_{\mathrm{unif}}$ saturating the same budget,
$\sigma_{\mathrm{unif}}^{2} = (\sum_c \Delta_c^{2})/(2\,\rho_{\mathrm{rel}})$,
the uniform-allocation objective is
\begin{equation}
J_{\mathrm{unif}}
 \;=\; \frac{1}{2\,\rho_{\mathrm{rel}}}
       \Bigl(\sum_{c}\Delta_c^{2}\Bigr)\Bigl(\sum_{c}s_c\Bigr).
\label{eq:unif-obj}
\end{equation}
By Cauchy--Schwarz $(\sum_c \Delta_c\sqrt{s_c})^{2} \le
(\sum_c \Delta_c^{2})(\sum_c s_c)$ with equality iff $s_c$ is
constant, so $J_{\mathrm{WF}}^{\star} \le J_{\mathrm{unif}}$. Strict
inequality holds whenever the $\{s_c\}$ are not all equal; with
$\Delta_c \equiv \Delta$ the ratio $\sigma^{\star}_c / \sigma^{\star}_{c'} =
(s_{c'}/s_c)^{1/4}$, bounded by $R^{1/4}$ where $R = s_{\max}/s_{\min}$.
This establishes Corollary~\ref{cor:r14} and Theorem~\ref{thm:wf}.\qed

% ---------------------------------------------------------------------
\section{Full proof of Theorem~\ref{thm:pate-sens} (PATE aggregate sensitivity)}
\label{app:pate-proof}

We restate. Let $\{\phi_T^{(k)}\}_{k=1}^K$ be teacher encoders trained
on disjoint patient cohorts $\mathcal{D}_k$, and let
$\mathrm{norm}(\cdot;\widehat{c}^{(k)})$ denote per-channel clipping
to cap $\widehat{c}^{(k)}$ followed by per-channel division by the
cap, so that each channel of $\mathrm{norm}(\phi_T^{(k)}(x);
\widehat{c}^{(k)})$ has $L_2$ norm $\le 1$.

\paragraph{Disjoint-cohort effect of a replace-one.}
Let $\mathcal{D}$ and $\mathcal{D}'$ differ by replacing exactly one
patient. Because cohorts are disjoint, the replaced patient belongs to
exactly one cohort, say $\mathcal{D}_{k^\star}$. Teacher $k^\star$ is
trained on $\mathcal{D}_{k^\star}$ in the original execution and on
$(\mathcal{D}_{k^\star})'$ in the neighbouring execution; teachers
$k \ne k^\star$ depend only on their own cohort and therefore satisfy
$\phi_T^{(k)} = \phi_T^{(k)\prime}$.

\paragraph{Per-channel bound on the aggregate.}
Let $\bar z(x;\mathcal{D}) = \tfrac{1}{K}\sum_{k=1}^{K}
\mathrm{norm}(\phi_T^{(k)}(x); \widehat{c}^{(k)})$ be the aggregate at
input $x$, where the $\phi_T^{(k)}$ depend implicitly on
$\mathcal{D}_k$. By the disjoint-cohort observation,
\begin{equation*}
\bar z(x;\mathcal{D}) - \bar z(x;\mathcal{D}')
 \;=\; \frac{1}{K}\,\bigl(
   \mathrm{norm}\bigl(\phi_T^{(k^\star)}(x);\widehat{c}^{(k^\star)}\bigr)
 \;-\;
   \mathrm{norm}\bigl(\phi_T^{(k^\star)\prime}(x);\widehat{c}^{(k^\star)}\bigr)
 \bigr).
\end{equation*}
Both normalised teacher outputs lie in the per-channel unit $L_2$ ball,
so their difference has per-channel $L_2$ norm $\le 2$ (achieved at
diametrically opposite points). Dividing by $K$:
\begin{equation*}
\bigl\|\, \bar z(x;\mathcal{D})_c - \bar z(x;\mathcal{D}')_c \,\bigr\|_2
 \;\le\; 2/K \quad \text{for every channel } c.
\end{equation*}
This is the per-channel $L_2$ replace-one sensitivity $\Delta_c$
required by the Gaussian-mechanism accounting in \eqref{eq:rho-rel},
so $\Delta_c = 2/K$ as claimed. \qed

\paragraph{Consequence.}
Substituting $\Delta_c = 2/K$ into the uniform allocation of
\eqref{eq:rho-rel} immediately gives \eqref{eq:pate-sigma}, i.e.,
$\sigma^{\mathrm{PATE}}(K) = (2/K)\sqrt{C/(2\rho_{\mathrm{rel}})}$,
showing that the noise scale decreases linearly in $K$ at fixed
budget. CANAL (Theorem~\ref{thm:wf}) plugs into this with $\Delta_c
\equiv 2/K$ replacing the $K\!=\!1$ sensitivity and is orthogonal to
the $K$ scaling: any channel allocation inherits the $K$-fold noise
reduction.

% ---------------------------------------------------------------------
%  BibTeX entries to add to references.bib (paste into your .bib):
%
%  @inproceedings{bun2016concentrated,
%    title={Concentrated Differential Privacy: Simplifications,
%           Extensions, and Lower Bounds},
%    author={Bun, Mark and Steinke, Thomas},
%    booktitle={Theory of Cryptography Conference (TCC)},
%    year={2016},
%  }
%
%  @inproceedings{lan2024gkd,
%    title={Gradient-Guided Knowledge Distillation for Object Detection},
%    author={Lan, Qizhen and Tian, Qing},
%    booktitle={IEEE/CVF Winter Conference on Applications of
%               Computer Vision (WACV)},
%    year={2024},
%  }
%
%  @inproceedings{papernot2017semi,
%    title={Semi-Supervised Knowledge Transfer for Deep Learning from
%           Private Training Data},
%    author={Papernot, Nicolas and Abadi, Mart\'in and Erlingsson,
%            \'Ulfar and Goodfellow, Ian and Talwar, Kunal},
%    booktitle={International Conference on Learning Representations
%               (ICLR)},
%    year={2017},
%  }
%
%  @inproceedings{papernot2018scalable,
%    title={Scalable Private Learning with PATE},
%    author={Papernot, Nicolas and Song, Shuang and Mironov, Ilya and
%            Raghunathan, Ananth and Talwar, Kunal and Erlingsson,
%            \'Ulfar},
%    booktitle={International Conference on Learning Representations
%               (ICLR)},
%    year={2018},
%  }
% ---------------------------------------------------------------------
